\begin{axiombox}{Axiom (Substitution of Equals)}
\begin{axiom}[Substitution of Equals]
\label{ax:equality-substitution}
For every formula $\varphi(z)$ whose substitution is admissible,
\[
x=y \Longrightarrow \bigl(\varphi(x)\Longleftrightarrow\varphi(y)\bigr).
\]
\end{axiom}
\end{axiombox}

\begin{remark*}[Standard quantified statement]
\[
\forall x\,\forall y\,
\Bigl(x=y \Longrightarrow
\bigl(\varphi(x)\Longleftrightarrow\varphi(y)\bigr)\Bigr).
\]
\end{remark*}

\begin{remark*}[Interpretation]
Substitution is the operational meaning of equality. Equal objects may replace
one another in legitimate statements without changing truth.
\end{remark*}

\begin{dependencies}
\begin{itemize}
  \item \hyperref[ax:leibniz-law]{Leibniz's Law}
\end{itemize}
\end{dependencies}

\begin{propositionbox}{Proposition (Substitution Preserves Predicates)}
\begin{proposition}[Substitution Preserves Predicates]
\label{prop:substitution-preserves-predicates}
\[
x=y \Longrightarrow \bigl(P(x)\Longleftrightarrow P(y)\bigr).
\]
\end{proposition}
\end{propositionbox}

\begin{remark*}[Standard quantified statement]
\[
\forall x\,\forall y,\,
x=y \Longrightarrow \bigl(P(x)\Longleftrightarrow P(y)\bigr).
\]
\end{remark*}

\begin{remark*}[Interpretation]
Predicates cannot distinguish equal arguments.
\end{remark*}

\begin{dependencies}
\begin{itemize}
  \item \hyperref[ax:equality-substitution]{Substitution of Equals}
\end{itemize}
\end{dependencies}

\begin{propositionbox}{Proposition (Substitution Preserves Relations on the Left)}
\begin{proposition}[Substitution Preserves Relations on the Left]
\label{prop:substitution-preserves-relations-left}
\[
x=y \Longrightarrow \bigl(R(x,z)\Longleftrightarrow R(y,z)\bigr).
\]
\end{proposition}
\end{propositionbox}

\begin{remark*}[Standard quantified statement]
\[
\forall x\,\forall y\,\forall z,\,
x=y \Longrightarrow
\bigl(R(x,z)\Longleftrightarrow R(y,z)\bigr).
\]
\end{remark*}

\begin{remark*}[Interpretation]
Relations respect equality in their first coordinate.
\end{remark*}

\begin{dependencies}
\begin{itemize}
  \item \hyperref[ax:equality-substitution]{Substitution of Equals}
\end{itemize}
\end{dependencies}

\begin{propositionbox}{Proposition (Substitution Preserves Relations on the Right)}
\begin{proposition}[Substitution Preserves Relations on the Right]
\label{prop:substitution-preserves-relations-right}
\[
x=y \Longrightarrow \bigl(R(z,x)\Longleftrightarrow R(z,y)\bigr).
\]
\end{proposition}
\end{propositionbox}

\begin{remark*}[Standard quantified statement]
\[
\forall x\,\forall y\,\forall z,\,
x=y \Longrightarrow
\bigl(R(z,x)\Longleftrightarrow R(z,y)\bigr).
\]
\end{remark*}

\begin{remark*}[Interpretation]
Relations respect equality in their second coordinate.
\end{remark*}

\begin{dependencies}
\begin{itemize}
  \item \hyperref[ax:equality-substitution]{Substitution of Equals}
\end{itemize}
\end{dependencies}

\begin{propositionbox}{Proposition (Substitution Preserves Relations)}
\begin{proposition}[Substitution Preserves Relations]
\label{prop:substitution-preserves-relations-full}
\[
x=x' \land y=y' \Longrightarrow
\bigl(R(x,y)\Longleftrightarrow R(x',y')\bigr).
\]
\end{proposition}
\end{propositionbox}

\begin{remark*}[Standard quantified statement]
\[
\forall x,x',y,y',\,
\bigl(x=x'\land y=y'\bigr)
\Longrightarrow
\bigl(R(x,y)\Longleftrightarrow R(x',y')\bigr).
\]
\end{remark*}

\begin{remark*}[Interpretation]
Relations are congruent with respect to equality in all displayed arguments.
\end{remark*}

\begin{dependencies}
\begin{itemize}
  \item \hyperref[prop:substitution-preserves-relations-left]{Substitution Preserves Relations on the Left}
  \item \hyperref[prop:substitution-preserves-relations-right]{Substitution Preserves Relations on the Right}
\end{itemize}
\end{dependencies}

\begin{propositionbox}{Proposition (Substitution Preserves Functions)}
\begin{proposition}[Substitution Preserves Functions]
\label{prop:substitution-preserves-functions}
\[
x=y \Longrightarrow f(x)=f(y).
\]
\end{proposition}
\end{propositionbox}

\begin{remark*}[Standard quantified statement]
\[
\forall x\,\forall y,\,
x=y \Longrightarrow f(x)=f(y).
\]
\end{remark*}

\begin{remark*}[Interpretation]
Functions send equal inputs to equal outputs.
\end{remark*}

\begin{dependencies}
\begin{itemize}
  \item \hyperref[ax:equality-substitution]{Substitution of Equals}
\end{itemize}
\end{dependencies}

\begin{propositionbox}{Proposition (Left Congruence for Operations)}
\begin{proposition}[Left Congruence for Operations]
\label{prop:substitution-preserves-operations-left}
\[
x=x' \Longrightarrow x\mathbin{O}y=x'\mathbin{O}y.
\]
\end{proposition}
\end{propositionbox}

\begin{remark*}[Standard quantified statement]
\[
\forall x,x',y,\,
x=x' \Longrightarrow x\mathbin{O}y=x'\mathbin{O}y.
\]
\end{remark*}

\begin{remark*}[Interpretation]
A binary operation respects equality in its left argument.
\end{remark*}

\begin{dependencies}
\begin{itemize}
  \item \hyperref[prop:substitution-preserves-functions]{Substitution Preserves Functions}
\end{itemize}
\end{dependencies}

\begin{propositionbox}{Proposition (Right Congruence for Operations)}
\begin{proposition}[Right Congruence for Operations]
\label{prop:substitution-preserves-operations-right}
\[
y=y' \Longrightarrow x\mathbin{O}y=x\mathbin{O}y'.
\]
\end{proposition}
\end{propositionbox}

\begin{remark*}[Standard quantified statement]
\[
\forall x,y,y',\,
y=y' \Longrightarrow x\mathbin{O}y=x\mathbin{O}y'.
\]
\end{remark*}

\begin{remark*}[Interpretation]
A binary operation respects equality in its right argument.
\end{remark*}

\begin{dependencies}
\begin{itemize}
  \item \hyperref[prop:substitution-preserves-functions]{Substitution Preserves Functions}
\end{itemize}
\end{dependencies}

\begin{propositionbox}{Proposition (Full Congruence for Operations)}
\begin{proposition}[Full Congruence for Operations]
\label{prop:substitution-preserves-operations-full}
\[
x=x' \land y=y' \Longrightarrow x\mathbin{O}y=x'\mathbin{O}y'.
\]
\end{proposition}
\end{propositionbox}

\begin{remark*}[Standard quantified statement]
\[
\forall x,x',y,y',\,
\bigl(x=x'\land y=y'\bigr)
\Longrightarrow
x\mathbin{O}y=x'\mathbin{O}y'.
\]
\end{remark*}

\begin{remark*}[Interpretation]
Binary operations are congruent with respect to equality in both arguments.
\end{remark*}

\begin{dependencies}
\begin{itemize}
  \item \hyperref[prop:substitution-preserves-operations-left]{Left Congruence for Operations}
  \item \hyperref[prop:substitution-preserves-operations-right]{Right Congruence for Operations}
\end{itemize}
\end{dependencies}

\begin{propositionbox}{Proposition (Congruence with Respect to Equality Is Automatic)}
\begin{proposition}[Congruence with Respect to Equality Is Automatic]
\label{prop:congruence-with-respect-to-equality-is-automatic}
Every function, predicate, relation, and operation respects equality in its
displayed arguments.
\end{proposition}
\end{propositionbox}

\begin{remark*}[Standard quantified statement]
\[
\begin{aligned}
x=y&\Longrightarrow f(x)=f(y),\\
x=y&\Longrightarrow \bigl(P(x)\Longleftrightarrow P(y)\bigr),\\
x=x'\land y=y'
&\Longrightarrow \bigl(R(x,y)\Longleftrightarrow R(x',y')\bigr),\\
x=x'\land y=y'
&\Longrightarrow x\mathbin{O}y=x'\mathbin{O}y'.
\end{aligned}
\]
\end{remark*}

\begin{remark*}[Interpretation]
Respecting equality is automatic. Respecting a nontrivial equivalence relation
is not automatic; that is why quotient constructions must prove separate
respect lemmas before operations, predicates, or relations descend to classes.
\end{remark*}

\begin{dependencies}
\begin{itemize}
  \item \hyperref[prop:substitution-preserves-predicates]{Substitution Preserves Predicates}
  \item \hyperref[prop:substitution-preserves-relations-full]{Substitution Preserves Relations}
  \item \hyperref[prop:substitution-preserves-functions]{Substitution Preserves Functions}
  \item \hyperref[prop:substitution-preserves-operations-full]{Full Congruence for Operations}
\end{itemize}
\end{dependencies}
